\section{Praktikum}
\label{sec:praktikum-bab9}

\begin{praktikumbox}
\vspace{1.1em}
\textbf{Judul:} Menggunakan layanan BIOS/DOS dan memeriksa tabel
vektor interupsi.\\[0.3em]
\textbf{Sub-CPMK:} 2.3 \quad\textbullet\quad
\textbf{Alokasi waktu:} 100 menit \quad\textbullet\quad
\textbf{Bentuk:} individu
\end{praktikumbox}

\subsection{Tujuan Praktikum}

Setelah menyelesaikan praktikum ini, mahasiswa mampu menjalankan
program dengan layanan \code{INT 21h}, \code{INT 10h}, dan
\code{INT 16h}, serta menjelaskan isi tabel vektor interupsi.

\subsection{Alat dan Bahan}

\begin{itemize}
    \item Komputer dengan emu8086 terpasang.
    \item DOSBox dengan TASM, TLINK, dan DEBUG.
    \item Berkas \file{input-nama.asm} dan \file{warna-layar.asm}
          dari \file{code/asm/}.
    \item Lembar jurnal praktikum.
\end{itemize}

\subsection{Langkah Kerja}

\begin{enumerate}
    \item \textbf{Menguji masukan dan keluaran DOS.}
    \begin{enumerate}
        \item Rakit \file{input-nama.asm} dengan TASM, lalu link
              dengan TLINK.
        \item Jalankan program dan catat pertanyaan serta keluaran
              yang muncul.
        \item Buka program dengan DEBUG, lalu periksa penyangga
              \code{nama} pada segmen data.
        \item Catat byte kapasitas, byte jumlah karakter, dan
              karakter yang tersimpan.
    \end{enumerate}

    \item \textbf{Menguji layanan video BIOS.}
    \begin{enumerate}
        \item Rakit dan jalankan \file{warna-layar.asm}.
        \item Catat posisi judul dan garis yang tampil.
        \item Ubah atribut warna judul dari \hex{0E} menjadi
              \hex{0A}, lalu rakit ulang dan amati perubahannya.
        \item Ubah posisi baris judul dari 5 menjadi 3 dan kolom dari
              10 menjadi 20. Catat hasilnya.
    \end{enumerate}

    \item \textbf{Memeriksa tabel vektor interupsi.}
    \begin{enumerate}
        \item Jalankan DEBUG.
        \item Tampilkan vektor \code{INT 21h} dengan perintah
              \lstasm|D 0000:0084 L 4|.
        \item Catat dua byte offset dan dua byte segmen yang tampil.
        \item Hitung alamat fisik handler dengan rumus
              $(\text{segmen} \times \hex{10}) + \text{offset}$.
        \item Ulangi untuk \code{INT 10h} pada alamat
              \hex{0000}:\hex{0040}.
    \end{enumerate}

    \item \textbf{Menguji interupsi keyboard BIOS.}
    \begin{enumerate}
        \item Tulis program berikut sebagai \file{baca-tombol.asm}.

\begin{lstlisting}[language=MASM, numbers=none]
.MODEL SMALL
.CODE
START:
    MOV AH, 00h     ; baca tombol, menunggu
    INT 16h         ; hasil di AH (scan) dan AL (ASCII)
    MOV AH, 4Ch
    INT 21h
END START
\end{lstlisting}

        \item Rakit, link, lalu jalankan program.
        \item Tekan beberapa tombol dan jelaskan mengapa program baru
              berakhir setelah tombol ditekan.
        \item Periksa isi \reg{AX} setelah \code{INT 16h} dengan
              \eng{debugger}.
    \end{enumerate}

    \item \textbf{Modifikasi program.}
    \begin{enumerate}
        \item Ubah pertanyaan pada \file{input-nama.asm} sesuai
              kebutuhan.
        \item Ubah teks dan warna pada \file{warna-layar.asm}.
        \item Rakit ulang dan pastikan perubahan terlihat.
    \end{enumerate}
\end{enumerate}

\subsection{Hasil yang Dikumpulkan}

\begin{enumerate}
    \item Berkas \file{input-nama.asm} dan \file{warna-layar.asm}
          versi modifikasi.
    \item Tangkapan layar keluaran program masukan.
    \item Tangkapan layar tampilan berwarna.
    \item Tabel vektor \code{INT 21h} dan \code{INT 10h} beserta
          alamat fisik handler.
    \item Berkas \file{baca-tombol.asm} dan penjelasan hasilnya.
\end{enumerate}

\begin{asesmenbox}
\vspace{1.1em}
\textbf{Asesmen praktikum 7 (bobot 5\% dari komponen praktikum).}

\begin{tabular}{@{}p{6.2cm}cp{4.6cm}@{}}
    \toprule
    \textbf{Indikator} & \textbf{Bobot} & \textbf{Kriteria} \\
    \midrule
    Program masukan/keluaran DOS berjalan & 30\% &
    Pertanyaan, masukan, dan sapaan tampil dengan benar \\
    \addlinespace
    Program layanan video berjalan & 25\% &
    Teks dan warna tampil sesuai register yang diisi \\
    \addlinespace
    Ketepatan membaca tabel vektor & 30\% &
    Offset, segmen, dan alamat fisik handler dihitung benar \\
    \addlinespace
    Analisis interupsi keyboard & 15\% &
    Perilaku menunggu dan isi register dijelaskan \\
    \bottomrule
\end{tabular}
\end{asesmenbox}
